% DEPRECATED: This is an older draft. The canonical manuscript is in:
% Full project/01_Science/manuscript/SOL.tex
% This file is kept for reference only.

\documentclass[12pt, a4paper]{article}
\usepackage[utf8]{inputenc}
\usepackage[T1]{fontenc}
\usepackage[portuguese]{babel}
\usepackage{graphicx}
\usepackage{geometry}
\usepackage{amsmath}
\usepackage{booktabs}
\usepackage{caption}
\usepackage{subcaption}
\usepackage{natbib}
\usepackage{setspace}
\usepackage{hyperref}

\geometry{a4paper, margin=2.5cm}
\onehalfspacing

\begin{document}

% TÍTULO E AUTORES
\title{Um Framework Multi-Critério para Localização de Sistemas Solares Comunitários em Angola: Estudo de Caso da Província de [Nome]}
\author{
    Nome do Autor 1\thanks{Instituição, email} \and 
    Nome do Autor 2\thanks{Instituição, email} \and 
    Nome do Autor 3\thanks{Instituição, email}
}
\date{\today}
\maketitle

% RESUMO E PALAVRAS-CHAVE
\begin{abstract}
    \textbf{Contexto:} [Escreva aqui o contexto em 2-3 linhas] \\
    \textbf{Objetivo:} [Objetivo claro em 1-2 linhas] \\
    \textbf{Metodologia:} [Descrição metodológica concisa] \\
    \textbf{Resultados:} [Principais resultados encontrados] \\
    \textbf{Conclusões:} [Conclusão principal e implicações] \\
    \textbf{Palavras-chave:} energia solar comunitária; sistemas de informação geográfica; análise multi-critério; planeamento energético; Angola.
\end{abstract}

\newpage

% ÍNDICE
\tableofcontents
\newpage

% SEÇÃO 1: INTRODUÇÃO
\section{Introdução}
\subsection{Contexto Global e Relevância do Problema}
\subsubsection{O Desafio do Acesso à Energia em África}
\subsubsection{Os Objetivos de Desenvolvimento Sustentável e a Energia}

\subsection{Contexto Nacional Angolano}
\subsubsection{Panorama Energético de Angola}
\subsubsection{Desafios Específicos da Eletrificação Rural}

\subsection{Revisão da Literatura}
\subsubsection{Sistemas Solares Comunitários: Experiências Internacionais}
\subsubsection{Métodos de Localização de Infraestruturas Energéticas}
\subsubsection{Lacunas na Literatura Específica para Angola}

\subsection{Objetivos do Estudo}
\subsubsection{Objetivo Principal}
\subsubsection{Objetivos Específicos}

\subsection{Estrutura do Artigo}

% SEÇÃO 2: ENQUADRAMENTO TEÓRICO
\section{Enquadramento Teórico}
\subsection{Conceito de Sistemas Solares Comunitários}
\subsubsection{Definição e Características}
\subsubsection{Modelos de Negócio e Governança}

\subsection{Análise Multi-Critério para Decisão Espacial}
\subsubsection{Fundamentos da Análise Multi-Critério}
\subsubsection{Aplicação em Planeamento Energético}

\subsection{Sistemas de Informação Geográfica para Energias Renováveis}
\subsubsection{Ferramentas e Técnicas Relevantes}

% SEÇÃO 3: METODOLOGIA
\section{Metodologia}
\subsection{Área de Estudo}
\subsubsection{Justificação da Seleção}
\subsubsection{Caracterização da Região}

\subsection{Framework Metodológico}
\subsubsection{Visão Geral do Processo}

\subsection{Critérios de Seleção e Fontes de Dados}
\subsubsection{Critérios Técnicos}
\begin{itemize}
    \item Irradiação Solar
    \item Topografia e Inclinação
    \item Proximidade de Infraestruturas
\end{itemize}

\subsubsection{Critérios Socioeconômicos}
\begin{itemize}
    \item Densidade Populacional
    \item Índices de Pobreza
    \item Atividades Econômicas Locais
\end{itemize}

\subsubsection{Critérios Ambientais}
\begin{itemize}
    \item Uso do Solo
    \item Áreas Protegidas
    \item Impacto Visual
\end{itemize}

\subsubsection{Fontes de Dados e Pré-processamento}

\subsection{Ponderação dos Critérios}
\subsubsection{Método de Determinação de Pesos}
\subsubsection{Matriz de Comparação Par a Par}

\subsection{Análise em Sistema de Informação Geográfica}
\subsubsection{Processamento Espacial}
\subsubsection{Análise de Sobreposição}

\subsection{Avaliação de Tecnologias}
\subsubsection{Opções Tecnológicas Consideradas}
\subsubsection{Parâmetros de Comparação}

\subsection{Limitações Metodológicas}

% SEÇÃO 4: RESULTADOS
\section{Resultados}
\subsection{Caracterização da Área de Estudo}
\subsubsection{Análise dos Dados Base}

\subsection{Mapas de Aptidão por Critério Individual}
\subsubsection{Mapa de Potencial Solar}
\subsubsection{Mapa de Densidade Populacional}
\subsubsection{Mapa de Acesso à Rede Elétrica}
\subsubsection{Mapa de Compatibilidade de Uso do Solo}

\subsection{Mapa de Aptidão Integrada}
\subsubsection{Zonas de Prioridade Alta, Média e Baixa}
\subsubsection{Análise das Áreas Prioritárias}

\subsection{Avaliação Comparativa de Tecnologias}
\subsubsection{Tabela Comparativa de Opções Tecnológicas}
\subsubsection{Recomendações por Tipo de Zona}

\subsection{Sensibilidade da Análise}
\subsubsection{Análise de Sensibilidade dos Pesos}

% SEÇÃO 5: DISCUSSÃO
\section{Discussão}
\subsection{Interpretação dos Resultados Principais}
\subsubsection{Significado das Zonas Prioritárias Identificadas}
\subsubsection{Fatores Determinantes na Seleção}

\subsection{Comparação com Estudos Similares}
\subsubsection{Semelhanças e Diferenças com Outros Contextos}
\subsubsection{Especificidades do Caso Angolano}

\subsection{Implicações Práticas}
\subsubsection{Para Políticas Públicas}
\subsubsection{Para Investidores e Projetistas}
\subsubsection{Para as Comunidades Locais}

\subsection{Limitações do Estudo}
\subsubsection{Limitações de Dados}
\subsubsection{Limitações Metodológicas}
\subsubsection{Limitações de Escopo}

\subsection{Sugestões para Investigação Futura}
\subsubsection{Expansão Geográfica}
\subsubsection{Refinamento Metodológico}
\subsubsection{Análises Complementares}

% SEÇÃO 6: CONCLUSÕES
\section{Conclusões}
\subsection{Resumo das Principais Descobertas}
\subsection{Contribuições do Estudo}
\subsubsection{Contribuições Teóricas}
\subsubsection{Contribuições Práticas}

\subsection{Recomendações Finais}
\subsubsection{Recomendações para Ação Imediata}
\subsubsection{Recomendações para Políticas de Longo Prazo}

\subsection{Reflexão Final}

% AGRADECIMENTOS
\section*{Agradecimentos}
Os autores agradecem... [Aqui vão os agradecimentos específicos]

% REFERÊNCIAS
\begin{thebibliography}{99}
\bibitem{} [Sua primeira referência]
\bibitem{} [Sua segunda referência]
\end{thebibliography}

% APÊNDICES
\appendix
\section{Dados e Metodologia Complementar}
\subsection{Descrição Detalhada dos Dados}
\subsection{Scripts de Processamento em SIG}
\subsection{Questionário de Especialistas (se aplicável)}

\end{document}